\documentclass[11pt]{letter}
\usepackage[utf8]{inputenc}
\usepackage[T1]{fontenc}
\usepackage{geometry}
\usepackage{parskip}
\usepackage[hyphens]{url}

\geometry{letterpaper, margin=1in}

\signature{Alex Welcing}
\address{Lupine Science\\alex@lupine.io}

\begin{document}

\begin{letter}{Editor\\\textit{Integrating Materials and Manufacturing Innovation}}

\opening{Dear Editor,}

I am pleased to submit our manuscript entitled ``Forty Years of Interatomic Potentials, One Error Geometry: Hyper-Ribbon Universality and Why Pooled Benchmarks Mislead'' for consideration as a Technical Article in \textit{Integrating Materials and Manufacturing Innovation} (IMMI). This is a solo-authored work by Alex Welcing (Lupine Science).

\textbf{What is new.} This paper presents the first demonstration that elastic constant prediction errors for 559 interatomic potentials universally occupy hyper-ribbon manifolds with effective dimensionality 1.0--2.3 out of 3 (median 1.09)---the signature of sloppy model universality applied to materials simulation. We discover a striking BCC/FCC crystal-structure dichotomy in prediction accuracy, and present the first application of random-effects meta-analysis and formal aggregation-bias detection to interatomic potential benchmarking, revealing extreme between-element heterogeneity ($I^2 = 98.6\%$) that standard pooled benchmarks obscure. We further extend the error-geometry analysis to three foundation MLIPs (MACE-MP-0, CHGNet, Orb-v3), showing that foundation-model errors are directionally structured as well, that Born-stability failures concentrate in BCC refractory metals, and that classical alignment patterns do not transfer across modeling paradigms.

\textbf{Why IMMI.} IMMI is an ideal venue for this work because it uniquely publishes at the intersection of materials informatics, computational methodology, and manufacturing-relevant property prediction---the precise nexus of our contributions. The journal's emphasis on integrative approaches aligns with our goal of bridging statistical methodology with actionable insights for materials simulation.

This manuscript has not been submitted to any other journal for simultaneous consideration. No preprint version has been published. All code and data are available open-source at \url{https://github.com/alexwelcing/lupine}. The standalone analysis engine, \texttt{atlas-distill}, is released under the MIT License.

\textbf{Suggested Reviewers:}
\begin{itemize}
    \item \textbf{Prof.\ G{\'a}bor Cs{\'a}nyi} (University of Cambridge, \texttt{@cam.ac.uk})---Leading developer of machine-learning interatomic potentials (GAP, MACE); expertise in MLIP benchmarking and uncertainty quantification.
    \item \textbf{Prof.\ James P.\ Sethna} (Cornell University, \texttt{@cornell.edu})---Pioneer of sloppy model theory and information geometry; author of the seminal work on model manifold dimensionality reduction in complex systems.
    \item \textbf{Prof.\ Dallas R.\ Trinkle} (University of Illinois Urbana-Champaign, \texttt{@illinois.edu})---Expert in interatomic potential development, atomistic simulation of mechanical properties, and computational materials science methodology.
\end{itemize}

\closing{Sincerely,}

\end{letter}

\end{document}
